Induction heads---attention heads that complete patterns of the form $[A][B]\ldots[A] \rightarrow B$---are considered fundamental to in-context learning in transformers.
Yet they have been primarily identified and studied using random repeated token sequences, an artificial setup far from the natural language these models were trained on.
We ask: do the same heads that perform induction on random data also do so on naturalistic text, or are there ``hidden'' induction mechanisms that only activate on training-distribution-like inputs?
We measure per-head induction scores on both random repeated sequences and \owt passages in \gpttwo (124M parameters), complemented by output-value (OV) copying analysis and causal ablation experiments.
We find that the rank ordering of heads is highly conserved across conditions (Spearman $\rho = 0.84$, top-5 overlap = 80\%), ruling out the existence of naturalistic-only induction heads in the standard sense.
However, induction scores are 10--40$\times$ lower on natural text, and late-layer heads (layers 9--11) show dramatically \emph{stronger} copying behavior on natural text than on random data---up to 29$\times$ higher OV copying scores.
Ablating the 21 random-detected induction heads increases natural-text loss by 76.6\%.
Our results show that standard random-sequence detection identifies the right heads but underestimates their natural-language capabilities: the same induction heads implement a richer, distribution-sensitive form of pattern completion on naturalistic data.
